\documentclass[12pt,a4paper]{article}
\usepackage[utf8]{inputenc}
\usepackage[T1]{fontenc}
\usepackage[spanish,es-tabla]{babel}
\usepackage{lmodern}\usepackage{titlesec}\usepackage{tikz}\usepackage{lastpage}\usepackage{colortbl}
\usepackage{geometry}
\geometry{left=2.65cm,right=2.65cm,top=2.75cm,bottom=2.7cm}
\usepackage{setspace}
\onehalfspacing
\usepackage{array,tabularx,longtable,enumitem}
\usepackage{amssymb}
\usepackage[hidelinks]{hyperref}
\usepackage{fancyhdr}
\usepackage{microtype}
\usepackage{graphicx}
\usepackage{xcolor}
\definecolor{azul}{HTML}{15324C}
\definecolor{acento}{HTML}{B58A55}
\definecolor{gris}{HTML}{64748B}
\definecolor{fondo}{HTML}{F3F6F8}
\definecolor{linea}{HTML}{D9E2E8}
\hypersetup{colorlinks=true,linkcolor=azul,urlcolor=azul}
\titleformat{\section}{\Large\bfseries\color{azul}}{\thesection}{0.85em}{}[\vspace{0.25em}{\color{acento}\titlerule[1pt]}]
\titleformat{\subsection}{\large\bfseries\color{azul}}{\thesubsection}{0.7em}{}
\titlespacing*{\section}{0pt}{2.2ex plus 1ex}{1.4ex}
\titlespacing*{\subsection}{0pt}{1.7ex plus .5ex}{.8ex}
\pagestyle{fancy}
\fancyhf{}
\fancyhead[L]{\small\sffamily\color{gris} UTEQ \quad / \quad INGENIERÍA DE REQUERIMIENTOS}
\fancyhead[R]{\small\sffamily\color{acento} GA-TL-02}
\fancyfoot[L]{\scriptsize\sffamily\color{gris} Entrevista estructurada y Design Thinking}
\fancyfoot[R]{\scriptsize\sffamily\color{gris} \thepage\ / \pageref{LastPage}}
\renewcommand{\headrulewidth}{0.35pt}
\renewcommand{\footrulewidth}{0pt}
\setlength{\headheight}{17pt}
\setlength{\parindent}{0pt}
\setlength{\parskip}{6pt}
\renewcommand{\arraystretch}{1.25}
\setlist[enumerate]{itemsep=5pt,topsep=5pt}
\newcommand{\pendiente}[1]{\par\medskip\noindent\begingroup\setlength{\fboxsep}{11pt}\colorbox{fondo}{\parbox{\dimexpr\linewidth-22pt\relax}{\textcolor{azul}{\bfseries EN DESARROLLO}\\[4pt]\small #1}}\endgroup\par\medskip}
\newcommand{\etiqueta}[1]{{\sffamily\footnotesize\bfseries\color{acento} #1}}
\begin{document}
\begin{titlepage}
\thispagestyle{empty}
\begin{tikzpicture}[remember picture,overlay]
\fill[azul] (current page.north west) rectangle ([yshift=-7.6cm]current page.north east);
\fill[acento] ([xshift=2.6cm,yshift=-7.6cm]current page.north west) rectangle ([xshift=7.1cm,yshift=-7.48cm]current page.north west);
\draw[linea,line width=.7pt] ([xshift=2.6cm,yshift=2.0cm]current page.south west)--([xshift=-2.6cm,yshift=2.0cm]current page.south east);
\end{tikzpicture}
\vspace*{0.75cm}
{\sffamily\color{white}\small\bfseries UNIVERSIDAD TÉCNICA ESTATAL DE QUEVEDO\par}
\vspace{0.18cm}
{\sffamily\color{white}\small Facultad de Ciencias de la Computación \quad | \quad Ingeniería de Software\par}
\vspace{1.05cm}
{\sffamily\color{white}\large\bfseries INGENIERÍA DE REQUERIMIENTOS\par}
\vspace{0.45cm}
{\sffamily\color{white}\fontsize{30}{36}\selectfont\bfseries GA-TL-02\par}
\vspace{0.16cm}
{\sffamily\color{white}\large Taller de entrevista estructurada\\ y Design Thinking aplicados al PFC\par}
\vspace{1.9cm}
\etiqueta{CASO DE ESTUDIO}\par
\vspace{0.3cm}
{\Large\bfseries\color{azul} Sistema de Gestión de Arriendos y Residencias Estudiantiles\par}
{\large\color{gris} En los alrededores de la UTEQ\par}
\vspace{1.05cm}
\etiqueta{EQUIPO RESPONSABLE DE LA ACTIVIDAD}\par
\vspace{0.2cm}
{\normalsize Córdova Carreño Mayra Lucila\\ Herrera Morán Gilmar Javier\\ Rivas Bodero Jordy Fabián\par}
\vspace{0.75cm}
\etiqueta{DOCENTE}\par
\vspace{0.2cm}
Gleiston Cicerón Guerrero Ulloa, PhD\par
\vfill
{\small\color{gris} Cuarto nivel \quad / \quad Unidad 2 \quad / \quad Actividad formativa\par}
{\small\color{gris} Fecha límite: 21 de noviembre de 2026\par}
\vspace{1cm}
\end{titlepage}
\thispagestyle{fancy}
{\Large\bfseries\color{azul} Contenido}\par\vspace{0.4cm}
\tableofcontents
\newpage
\section{Presentación y alcance}
Este informe corresponde a la actividad formativa GA-TL-02 - FORM. El equipo de esta actividad está integrado por Córdova Carreño Mayra Lucila, Herrera Morán Gilmar Javier y Rivas Bodero Jordy Fabián. Para desarrollar el taller, el equipo seleccionó el proyecto \emph{Sistema de Gestión de Arriendos y Residencias Estudiantiles en los alrededores de la UTEQ}, que pertenece originalmente a otro equipo del Proyecto de Fin de Curso (PFC) en el que participa una integrante. Su elección como caso de estudio no cambia la autoría ni la conformación del equipo original del PFC.

El ejercicio se enfoca en el rol de \textbf{propietario de una residencia estudiantil}. Se pretende explorar sus prácticas de administración, sus dificultades y las necesidades que podrían orientar una futura elicitación de requisitos. Las situaciones descritas son hipótesis iniciales del caso, no resultados comprobados de una entrevista.

\section{Objetivo}
Diseñar y ensayar una entrevista estructurada dirigida al propietario de una residencia estudiantil, para identificar necesidades y dificultades relacionadas con la administración de arriendos y, a partir de la información obtenida, elaborar un mapa de empatía, definir un problema centrado en el usuario y corregir las preguntas que presenten dificultades.

\section{Organización del equipo}
\begin{tabularx}{\linewidth}{|p{3.0cm}|X|}
\hline
\textbf{Integrante} & \textbf{Responsabilidad principal} \\ \hline
Mayra Córdova & Guion de entrevista y consentimiento informado; entrevistadora durante el ensayo. \\ \hline
Jordy Rivas & Observación, registro, revisión de preguntas y manifiesto de evidencias. \\ \hline
Gilmar Herrera & Mapa de empatía y definición del problema; representación del propietario durante el ensayo. \\ \hline
\end{tabularx}

\subsection{Mayra Córdova: guion y entrevista}
Responsable del guion de entrevista, el consentimiento informado y la conducción de la entrevista.

Actividades asignadas:
\begin{itemize}[itemsep=3pt,topsep=3pt]
\item Elaborar el guion de entrevista estructurada (quince preguntas organizadas en cinco bloques).
\item Preparar el protocolo y el formulario de consentimiento informado.
\item Solicitar el consentimiento y la autorización de grabación de audio antes de iniciar.
\item Actuar como entrevistadora durante el ensayo.
\item Participar en la revisión del documento final.
\end{itemize}

\subsection{Jordy Rivas: observación, registro y evidencias}
Responsable de la observación, el registro de respuestas, la revisión de preguntas y el manifiesto de evidencias.

Actividades asignadas:
\begin{itemize}[itemsep=3pt,topsep=3pt]
\item Preparar el formato de observación y registro de respuestas.
\item Actuar como observador y registrador durante el ensayo, sin intervenir en el diálogo.
\item Registrar tiempos, síntesis de respuestas y dificultades detectadas.
\item Identificar las preguntas defectuosas y proponer su corrección con justificación.
\item Elaborar el manifiesto de evidencias (archivo, fecha, duración, enlace y huella SHA-256) y el archivo \texttt{README.md}.
\item Participar en la revisión del documento final.
\end{itemize}

\subsection{Gilmar Herrera: Design Thinking y definición del problema}
Responsable de Design Thinking y de la definición del problema.

Actividades asignadas:
\begin{itemize}[itemsep=3pt,topsep=3pt]
\item Preparar la plantilla del mapa de empatía.
\item Interpretar al propietario durante el ensayo.
\item Analizar las respuestas obtenidas en el ensayo.
\item Identificar necesidades y dificultades del usuario representado.
\item Completar el mapa de empatía con base en la información registrada.
\item Formular la definición del problema centrada en el usuario mediante «¿Cómo podríamos\ldots?».
\item Participar en la revisión del documento final.
\end{itemize}

\subsection{Responsabilidades compartidas}
Los tres integrantes deberán:
\begin{itemize}[itemsep=3pt,topsep=3pt]
\item Revisar el guion antes del ensayo.
\item Participar en el juego de roles.
\item Verificar el consentimiento y la grabación de audio.
\item Revisar la coherencia del mapa de empatía.
\item Comprobar que los resultados correspondan a las evidencias.
\item Integrar y revisar el documento \LaTeX.
\item Verificar la compilación del PDF.
\item Comprobar la accesibilidad del repositorio.
\item Verificar la entrega antes del cierre oficial.
\end{itemize}

\section{Guion de entrevista estructurada}
\subsection{Ficha de planificación}
\begin{tabularx}{\linewidth}{@{}p{4.4cm}X@{}}
Parte interesada: & Propietario de residencia estudiantil (representado en juego de roles).\\
Entrevistadora: & Córdova Carreño Mayra Lucila.\\
Entrevistado (rol): & Herrera Morán Gilmar Javier.\\
Observador: & Rivas Bodero Jordy Fabián.\\
Duración prevista: & Aproximadamente 10 minutos.\\
Modalidad: & Ensayo académico con grabación únicamente de audio, previa autorización.\\
Estado: & Guion preparado; ensayo realizado (ver sección de observaciones).\\
\end{tabularx}

\subsection{Apertura}
\begin{quote}\itshape\color{azul}
Buenos días o buenas tardes. Somos estudiantes de Ingeniería de Software de la Universidad Técnica Estatal de Quevedo. Estamos realizando una actividad académica para conocer cómo se administran los arriendos y las residencias estudiantiles, así como las dificultades y necesidades de sus propietarios.

La entrevista durará aproximadamente diez minutos. Su participación es voluntaria y puede decidir no responder alguna pregunta o finalizar su participación. Antes de comenzar, revisaremos el consentimiento informado y solicitaremos autorización específica para grabar únicamente audio.
\end{quote}

\subsection{Bloque I. Contexto y administración}
\begin{enumerate}[label=\textbf{P\arabic*.},leftmargin=1.5cm]
\item ¿Cómo comenzó a administrar habitaciones o residencias destinadas a estudiantes?
\item ¿Qué actividades realiza habitualmente para gestionar su residencia?
\item ¿Cómo lleva el registro de las habitaciones que están ocupadas y disponibles?
\end{enumerate}

\subsection{Bloque II. Contratos y arrendatarios}
\begin{enumerate}[label=\textbf{P\arabic*.},start=4,leftmargin=1.5cm]
\item ¿Cómo realiza el proceso de ingreso de un nuevo estudiante a la residencia?
\item ¿Qué datos solicita a un estudiante cuando inicia un arrendamiento?
\item ¿Cómo organiza y conserva los acuerdos o contratos de arrendamiento?
\end{enumerate}

\subsection{Bloque III. Pagos y servicios básicos}
\begin{enumerate}[label=\textbf{P\arabic*.},start=7,leftmargin=1.5cm]
\item ¿Cómo lleva el control de las mensualidades que deben pagar los estudiantes?
\item ¿Qué procedimiento sigue cuando un estudiante se retrasa en el pago?
\item ¿Cómo administra el cobro de los servicios básicos utilizados en la residencia?
\end{enumerate}

\subsection{Bloque IV. Comunicación y mantenimiento}
\begin{enumerate}[label=\textbf{P\arabic*.},start=10,leftmargin=1.5cm]
\item ¿Qué medios utiliza para comunicarse con los estudiantes que alquilan sus habitaciones?
\item ¿Cómo recibe y atiende los reportes de daños o problemas de mantenimiento?
\item ¿Qué dificultades ha experimentado al coordinar reparaciones o solicitudes de los estudiantes?
\end{enumerate}

\subsection{Bloque V. Necesidades y oportunidades de mejora}
\begin{enumerate}[label=\textbf{P\arabic*.},start=13,leftmargin=1.5cm]
\item ¿Cuáles son las principales dificultades que enfrenta al administrar los arriendos?
\item ¿Qué aspectos de la administración le gustaría mejorar?
\item ¿Qué información le resulta más difícil encontrar cuando necesita tomar decisiones sobre sus arriendos?
\end{enumerate}

\subsection{Cierre}
\begin{quote}\itshape\color{azul}
Muchas gracias por compartir su experiencia. La información proporcionada nos ayudará a comprender mejor las necesidades relacionadas con la administración de residencias estudiantiles. Antes de finalizar, ¿desea agregar algún comentario que considere importante?
\end{quote}
\textit{Nota:} La invitación del cierre no forma parte de las quince preguntas numeradas.

\section{Protocolo de consentimiento informado}
\subsection{Propósito y procedimiento}
Se invita a participar voluntariamente en un ensayo académico de entrevista estructurada para identificar necesidades, dificultades y expectativas vinculadas con la administración de residencias estudiantiles. El participante representará el rol de propietario. La duración prevista es de diez minutos.

\subsection{Voluntariedad y derechos}
La participación es voluntaria. El participante puede omitir preguntas o retirarse en cualquier momento, sin consecuencias académicas o personales. Antes de iniciar, podrá solicitar aclaraciones sobre el ejercicio.

\subsection{Registro de audio y confidencialidad}
Se solicitará autorización diferenciada para la grabación \textbf{solo de audio}. No se realizará video. El audio y cualquier consentimiento firmado se conservarán en un espacio institucional privado, compartido únicamente con las personas autorizadas para revisión académica. No se publicarán en GitHub. En el repositorio público se documentarán únicamente los metadatos y la huella SHA-256 de la evidencia, sin datos personales sensibles.

\subsection{Constancia para el ensayo (formulario de uso privado)}
\textbf{Participación:} $\square$ Acepto participar \quad $\square$ No acepto participar.

\textbf{Grabación de audio:} $\square$ Autorizo \quad $\square$ No autorizo.

\textbf{Nombre del participante:} \rule{8cm}{0.4pt}

\textbf{Firma del participante:} \rule{8cm}{0.4pt}

\textbf{Fecha:} \rule{5cm}{0.4pt}

\textbf{Entrevistadora:} \rule{8cm}{0.4pt}

\textit{Este formulario es una plantilla sin completar. La copia firmada, si se utiliza, debe permanecer fuera del repositorio público. Si no se autoriza el audio, no se graba y deberá resolverse cómo cumplir el requisito de evidencia de la actividad.}

\section{Ensayo y registro de observaciones}
\subsection{Rol del observador y registrador}
Jordy Rivas actúa como observador y registrador durante el ensayo. No interviene en el diálogo y se limita a:
\begin{itemize}[itemsep=3pt,topsep=3pt]
\item Verificar antes de iniciar que el consentimiento esté revisado y que la autorización de audio esté decidida.
\item Controlar el tiempo de cada fase y anotar la hora real de inicio y fin.
\item Registrar para cada pregunta (P1--P15) una síntesis de la respuesta y si presentó algún defecto (ambigüedad, inducción, extensión excesiva u otro).
\item Anotar dificultades operativas, de comprensión, de neutralidad y técnicas (audio).
\item Conservar el audio en el espacio institucional privado y registrar su huella SHA-256 en el manifiesto.
\end{itemize}

\subsection{Formato de observación del ensayo}
{\footnotesize
\renewcommand{\arraystretch}{1.1}
\begin{tabularx}{\linewidth}{|p{3.0cm}|p{1.9cm}|p{1.8cm}|X|}
\hline
\textbf{Tramo / Fase} & \textbf{Previsto} & \textbf{Real} & \textbf{Observaciones de cumplimiento} \\ \hline
Apertura y consentimiento & Min 0--1 & 01:05 & Presentación y autorización de solo audio. \\ \hline
Bloques I--V (P1--P15) & Min 1--7 & 06:20 & Recorrido completo con respuestas detalladas. \\ \hline
Aclaraciones & Min 7--9 & 01:45 & Una pregunta opcional de profundización sobre el cálculo de servicios (tras P9). \\ \hline
Cierre & Min 9--10 & 00:50 & Cierre cordial antes de detener la grabación. \\ \hline
\end{tabularx}
}

\vspace{0.2cm}

{\scriptsize
\renewcommand{\arraystretch}{1.05}
\begin{longtable}{|p{1.0cm}|p{4.6cm}|p{1.7cm}|p{5.8cm}|}
\hline
\textbf{ID} & \textbf{Síntesis de respuesta} & \textbf{¿Defecto?} & \textbf{Registro de observación} \\ \hline
\endfirsthead
\hline
\textbf{ID} & \textbf{Síntesis de respuesta} & \textbf{¿Defecto?} & \textbf{Registro de observación} \\ \hline
\endhead
P1 & Comenzó ofreciendo habitaciones disponibles cerca de la universidad, por la demanda de estudiantes. & $\square$ No & Clara; aporta contexto amplio al inicio. \\ \hline
P2 & Revisa pagos y disponibilidad, atiende consultas y coordina limpieza o reparaciones. & $\square$ No & Fluida; reúne varias rutinas en una sola respuesta. \\ \hline
P3 & Anota la ocupación en un cuaderno y revisa mensajes de WhatsApp; compara ambos registros. & $\square$ No & Clara; evidencia duplicidad de registros. \\ \hline
P4 & Consulta de disponibilidad, visita, conversación sobre precio y condiciones, acuerdo e ingreso. & $\square$ No & Buena fluidez; el acuerdo parece informal. \\ \hline
P5 & Pide nombre completo, identificación, contacto y datos básicos. & $\square$ No & Respuesta directa y concreta. \\ \hline
P6 & Guarda documentos físicos en una carpeta; parte de la información queda en mensajes del teléfono. & $\square$ No & Evidencia dispersión entre papel y mensajería. \\ \hline
P7 & Anota fechas y valores en una libreta y revisa comprobantes de transferencias. & $\square$ No & Fluida; exige contrastar más de una fuente. \\ \hline
P8 & Escribe para recordar la fecha acordada y conversa para fijar cuándo pagará. & $\blacksquare$ Sí & Natural; seguimiento flexible; la pregunta asume un procedimiento fijo. \\ \hline
P9 & Depende del servicio: unos gastos son compartidos y otros incluidos; calcula lo que corresponde a cada persona. & $\blacksquare$ Sí & Mayor complejidad; la pregunta agrupa varios servicios y el cálculo. \\ \hline
P9+ & (Profundización) Debe buscar recibos y anotaciones; toma tiempo cuando los valores cambian. & $\square$ No & Pregunta opcional de aclaración; aporta detalle sobre el esfuerzo de revisión. \\ \hline
P10 & Usa WhatsApp, llamadas y trato presencial. & $\square$ No & Respuesta rápida y clara sobre canales. \\ \hline
P11 & Recibe avisos por mensaje o directamente; revisa y repara o contacta a alguien que pueda hacerlo. & $\square$ No & Adecuada y clara. \\ \hline
P12 & Es difícil coincidir en horario con técnico y estudiante; hay retrasos por falta de materiales o de disponibilidad del encargado. & $\blacksquare$ Sí & Detalla cuellos de botella; la pregunta presupone dificultades. \\ \hline
P13 & Mantener organizada la información: pagos, contratos, habitaciones y mantenimiento están en distintos lugares. & $\square$ No & Resume el principal problema operativo. \\ \hline
P14 & Organizar mejor pagos y contratos, tener claridad de disponibilidad y consultar reportes de mantenimiento. & $\square$ No & Enfocada en orden y trazabilidad. \\ \hline
P15 & Cuesta saber rápido cuánto está pendiente de pago, qué habitaciones estarán libres y cuánto se gastó en servicios o reparaciones. & $\blacksquare$ Sí & Útil para detectar vacíos de información; la pregunta es amplia. \\ \hline
Cierre & Sería útil tener la información organizada en un solo lugar. & $\square$ No & Comentario espontáneo en la invitación de cierre. \\ \hline
\end{longtable}
}

\subsection{Registro de dificultades detectadas}
{\footnotesize
\renewcommand{\arraystretch}{1.1}
\begin{tabularx}{\linewidth}{|p{3.1cm}|p{2.0cm}|X|p{1.8cm}|}
\hline
\textbf{Tipo} & \textbf{ID / Fase} & \textbf{Descripción de la dificultad} & \textbf{Estado} \\ \hline
Operativa / tiempos & P3, P6, P7, P9+, P12 & Verificar pagos en varios medios, recalcular servicios y cuadrar horarios entre técnico y estudiante alarga tiempos. & Registrada \\ \hline
Comprensión & P9, P15 & Términos amplios o agrupados admiten distintas lecturas sin delimitar el alcance. & Registrada \\ \hline
Neutralidad & P8, P12 & Algunas preguntas asumen rutinas fijas o dificultades previas. & Registrada \\ \hline
Técnica / audio & Apertura/Cierre & Verificación de nitidez de voz y grabación exclusiva de audio con consentimiento. & Registrada \\ \hline
\end{tabularx}
}

\section{Mapa de empatía (Design Thinking)}
\subsection{Base de información}
El mapa se construye solo con la síntesis de respuestas registrada en la tabla de observaciones (P1--P15). Cada punto indica entre paréntesis la pregunta de donde proviene. Los cuadrantes \textbf{Dice} y \textbf{Hace} recogen lo expresado o descrito por el participante; los cuadrantes \textbf{Piensa} y \textbf{Siente}, así como el análisis de necesidades, son \textbf{inferencias del equipo} y deben validarse con el audio y con entrevistas reales a propietarios.

\subsection{Mapa de empatía del propietario}
{\footnotesize
\renewcommand{\arraystretch}{1.1}
\begin{tabularx}{\linewidth}{|X|X|}
\hline
\rowcolor{fondo}\textbf{\color{azul}DICE} (expresado) & \textbf{\color{azul}HACE} (descrito) \\ \hline
\begin{itemize}[leftmargin=*,itemsep=1pt,topsep=2pt,label=--]
\item {Su principal dificultad es mantener organizada la información: pagos, contratos, habitaciones y mantenimiento están en distintos lugares (P13).}
\item {Quiere organizar mejor pagos y contratos, ver la disponibilidad y consultar reportes de mantenimiento (P14).}
\item {Le cuesta saber rápido cuánto está pendiente, qué habitaciones estarán libres y cuánto gastó en servicios y reparaciones (P15).}
\item {Es complicado coordinar el horario del técnico y del estudiante, y hay retrasos por falta de materiales o de disponibilidad del encargado (P12).}
\item {Sería útil tener la información organizada en un solo lugar (cierre).}
\end{itemize}
&
\begin{itemize}[leftmargin=*,itemsep=1pt,topsep=2pt,label=--]
\item {Anota la ocupación en un cuaderno y los pagos en una libreta; revisa mensajes y comprobantes (P3, P7).}
\item {Guarda documentos físicos en una carpeta; parte de la información queda en mensajes (P6).}
\item {Escribe para recordar la fecha acordada y conversa para fijar cuándo pagará (P8).}
\item {Calcula cuánto corresponde a cada persona en los consumos compartidos (P9).}
\item {Se comunica por WhatsApp, llamadas y en persona; repara o contacta a quien pueda hacerlo (P10, P11).}
\item {Pide nombre completo, identificación, contacto y datos básicos al ingresar un estudiante (P5).}
\end{itemize} \\ \hline
\rowcolor{fondo}\textbf{\color{azul}PIENSA} (inferencia) & \textbf{\color{azul}SIENTE} (inferencia) \\ \hline
\begin{itemize}[leftmargin=*,itemsep=1pt,topsep=2pt,label=--]
\item {Que tener cobros, contratos y disponibilidad ordenados facilitaría su administración (P14).}
\item {Que conversar con los estudiantes ante un atraso forma parte de su manera de administrar (P8).}
\item {Que la rapidez de las reparaciones depende en parte de factores ajenos a él (P12).}
\end{itemize}
&
\begin{itemize}[leftmargin=*,itemsep=1pt,topsep=2pt,label=--]
\item {Carga al tener que revisar varios lugares para confirmar un acuerdo o un pago (P6, P7).}
\item {Incertidumbre sobre cuánto está pendiente y cuánto se ha gastado (P15).}
\item {Desgaste por revisar recibos y anotaciones y por los retrasos en reparaciones (P9+, P12).}
\end{itemize} \\ \hline
\end{tabularx}
}

\subsection{Dolores y ganancias}
{\footnotesize
\renewcommand{\arraystretch}{1.1}
\begin{tabularx}{\linewidth}{|X|X|}
\hline
\rowcolor{fondo}\textbf{\color{azul}DOLORES} & \textbf{\color{azul}GANANCIAS / NECESIDADES} \\ \hline
\begin{itemize}[leftmargin=*,itemsep=1pt,topsep=2pt,label=--]
\item {Información de pagos, contratos y ocupación dispersa entre cuaderno, papel y mensajería (P3, P6, P7, P13).}
\item {Cálculo de servicios compartidos y su revisión posterior con recibos y anotaciones (P9, P9+).}
\item {Registros de ocupación duplicados (cuaderno y WhatsApp) que hay que comparar (P3).}
\item {Reparaciones condicionadas por el horario del técnico y del estudiante y por la falta de materiales (P12).}
\item {Seguimiento de pagos atrasados que depende de mensajes manuales (P8).}
\end{itemize}
&
\begin{itemize}[leftmargin=*,itemsep=1pt,topsep=2pt,label=--]
\item {Organizar pagos y contratos y ver la disponibilidad de habitaciones (P14).}
\item {Contar con trazabilidad de acuerdos y pagos (P4, P6, P7).}
\item {Saber con claridad cuánto está pendiente, qué habitaciones estarán libres y cuánto se gastó (P15).}
\item {Calcular y revisar el cobro de servicios compartidos con rapidez (P9, P9+).}
\item {Consultar fácilmente los reportes de mantenimiento (P14).}
\end{itemize} \\ \hline
\end{tabularx}
}


\section{Definición del problema centrada en el usuario}
\subsection{Planteamiento}
\textbf{Usuario:} propietario de una residencia estudiantil cerca de la UTEQ que administra los arriendos con cuaderno, comprobantes y mensajes.

\textbf{Necesidad:} saber con claridad el estado de pagos, ocupación y gastos, y tener ordenados los acuerdos con los estudiantes.

\textbf{Hallazgo:} la información existe, pero está repartida en varios medios, por lo que confirmar un pago, un acuerdo o el cálculo de servicios exige revisar varios lugares y consume tiempo (P3, P6, P7, P9+, P15).

\subsection{Pregunta «¿Cómo podríamos\ldots?»}
\begin{quote}\itshape\color{azul}
¿Cómo podríamos ayudar a los propietarios de residencias estudiantiles a conocer con claridad el estado de los pagos, la ocupación y los gastos de sus arriendos, aunque hoy esa información se lleve en varios medios?
\end{quote}
La pregunta no impone una solución tecnológica: describe el resultado que el propietario necesita y deja abiertas las alternativas.

\subsection{Preguntas complementarias}
\begin{itemize}[itemsep=3pt,topsep=3pt]
\item ¿Cómo podríamos facilitar el seguimiento de los pagos atrasados sin perder el trato flexible que el propietario tiene con los estudiantes?
\item ¿Cómo podríamos hacer más transparente el cálculo y el cobro de los servicios básicos compartidos?
\item ¿Cómo podríamos ayudar a coordinar las reparaciones considerando la disponibilidad de técnicos, estudiantes y materiales?
\end{itemize}

\textit{Nota:} la definición parte de un ensayo con juego de roles, por lo que es una hipótesis que debe confirmarse con propietarios reales.

\section{Revisión y ajustes de preguntas}
\subsection{Propuesta de correcciones y justificación}
{\footnotesize
\renewcommand{\arraystretch}{1.1}
\begin{tabularx}{\linewidth}{|p{0.9cm}|p{2.7cm}|p{2.2cm}|p{3.0cm}|X|}
\hline
\textbf{ID} & \textbf{Original} & \textbf{Defecto} & \textbf{Versión corregida} & \textbf{Justificación} \\ \hline
P8 & ¿Qué procedimiento sigue cuando un estudiante se retrasa en el pago? & Inductiva (asume protocolo fijo) & ¿Cómo suele actuar cuando un pago mensual no se realiza en la fecha acordada? & Evita imponer un procedimiento rígido y recoge prácticas reales. \\ \hline
P9 & ¿Cómo administra el cobro de los servicios básicos utilizados en la residencia? & Ambigua (agrupa rubros y cálculo) & ¿Cómo calcula y cobra los gastos de agua, electricidad e internet a los estudiantes? & Precisa servicios del contexto y separa cálculo de cobro. \\ \hline
P12 & ¿Qué dificultades ha experimentado al coordinar reparaciones o solicitudes de los estudiantes? & Inductiva (presupone fallas) & ¿Qué ocurre cuando necesita coordinar una reparación o atención solicitada por un estudiante? & Plantea la consulta sin sesgo para conocer proceso y obstáculos. \\ \hline
P15 & ¿Qué información le resulta más difícil encontrar cuando necesita tomar decisiones sobre sus arriendos? & Ambigua (decisiones sin acotar) & ¿Qué datos necesita consultar con mayor frecuencia al revisar pagos, ocupación o gastos? & Delimita el contexto de consulta para obtener respuestas concretas. \\ \hline
\end{tabularx}
}

\section{Evidencias y manifiesto}
\subsection{Manifiesto de evidencias}
{\footnotesize
\renewcommand{\arraystretch}{1.1}
\begin{tabularx}{\linewidth}{|p{3.6cm}|p{2.0cm}|p{1.6cm}|X|}
\hline
\textbf{Archivo} & \textbf{Fecha} & \textbf{Duración} & \textbf{Link} \\ \hline
Audio del ensayo (nombre por definir) & \textit{por completar} & \textit{por completar} & \textit{Enlace al almacenamiento institucional privado (por completar)} \\ \hline
\end{tabularx}
}

\textbf{Huella SHA-256 del audio:} \textit{por completar tras el ensayo} (se genera con \texttt{sha256sum archivo}). En el repositorio público solo se publican estos metadatos y la huella; el audio y el consentimiento firmado no se suben a GitHub.

\subsection{Archivo README del repositorio}
El archivo \texttt{README.md} de la carpeta \texttt{GA/TL-02} describe el contenido de la entrega, el equipo, la compilación y las reglas de privacidad de las evidencias.

\section{Procedimiento de compilación y revisión del documento final}
\subsection{Procedimiento de compilación}
\begin{quote}\ttfamily\small
git clone https://github.com/mcordovaac2/ISR-401-CHR.git\\
cd ISR-401-CHR/GA/TL-02\\
pdflatex -interaction=nonstopmode -halt-on-error main.tex\\
bibtex main\\
pdflatex -interaction=nonstopmode -halt-on-error main.tex\\
pdflatex -interaction=nonstopmode -halt-on-error main.tex
\end{quote}

\subsection{Revisión del documento final}
Registro de la revisión realizada por Jordy Rivas:
{\footnotesize
\renewcommand{\arraystretch}{1.1}
\begin{tabularx}{\linewidth}{|p{4.6cm}|X|}
\hline
\rowcolor{fondo}\textbf{Aspecto revisado} & \textbf{Resultado} \\ \hline
Observaciones y ajustes de preguntas & Las preguntas marcadas con defecto (P8, P9, P12 y P15) coinciden con la tabla de correcciones. \\ \hline
Mapa de empatía y problema & Cada punto del mapa cita la pregunta de origen (P1--P15); la pregunta «¿Cómo podríamos\ldots?» se deriva de esos hallazgos. \\ \hline
Manifiesto de evidencias & Estructura lista; nombre del archivo de audio, fecha, duración, enlace y huella SHA-256 se completan con los datos del ensayo. \\ \hline
Compilación del PDF & Compila sin errores con el procedimiento indicado. \\ \hline
Entrega & Fecha límite: 21 de noviembre de 2026. \\ \hline
\end{tabularx}
}


\subsection{Lista de verificación del equipo}
{\footnotesize
\begin{tabularx}{\linewidth}{|X|c|c|c|}
\hline
\textbf{Actividad} & \textbf{Mayra} & \textbf{Gilmar} & \textbf{Jordy} \\ \hline
Participó en el juego de roles de 10 minutos & $\square$ & $\square$ & $\square$ \\ \hline
Revisó el documento final & $\square$ & $\square$ & $\square$ \\ \hline
Verificó el manifiesto de evidencias y la huella SHA-256 & $\square$ & $\square$ & $\square$ \\ \hline
Confirmó la compilación limpia del PDF & $\square$ & $\square$ & $\square$ \\ \hline
\end{tabularx}
}

\section{Declaración de apoyo con inteligencia artificial}
Se utilizaron herramientas de inteligencia artificial como apoyo: \textbf{ChatGPT (OpenAI, modelo GPT-6)} para la redacción inicial, la revisión de neutralidad de preguntas y la preparación de la estructura \LaTeX, y \textbf{Claude (Anthropic, modelo Claude Sonnet 5.5)} para redactar la descripción de roles, el rol del observador, el borrador del mapa de empatía y de la definición del problema, el archivo \texttt{README.md} y la compilación del documento. La verificación de la pertinencia de las preguntas, el consentimiento, el ensayo, la fidelidad de los resultados y la compilación corresponde a los integrantes del equipo.

\section{Referencias}
\textit{Pendiente: incorporar únicamente referencias realmente consultadas y verificadas, con formato IEEE. No se incluyen entradas bibliográficas ficticias.}


% ==========================================================
% ANEXO A - CONSENTIMIENTO INFORMADO
% ==========================================================

\clearpage
\appendix

\section{Formulario de consentimiento informado}

\begin{center}
    {\large\bfseries CONSENTIMIENTO INFORMADO}\\[0.3cm]
    {\small Entrevista académica mediante juego de roles}\\[0.2cm]
    {\small GA-TL-02 - FORM | ISR-401}
\end{center}

\vspace{0.4cm}

\noindent
\textbf{Universidad:} Universidad Técnica Estatal de Quevedo\\
\textbf{Carrera:} Ingeniería de Software\\
\textbf{Asignatura:} Ingeniería de Requerimientos\\
\textbf{Proyecto:} Sistema de Gestión de Arriendos y
Residencias Estudiantiles en los alrededores de la UTEQ.

\subsection*{A.1. Propósito de la entrevista}

La presente entrevista forma parte de una actividad académica
orientada a identificar necesidades, dificultades y expectativas
relacionadas con la administración de arriendos y residencias
estudiantiles.

\subsection*{A.2. Procedimiento}

La actividad consiste en participar en una entrevista estructurada
mediante un ejercicio de juego de roles, en el que se representará
a una parte interesada del proyecto.

La duración estimada es de diez minutos.

\subsection*{A.3. Participación voluntaria}

La participación es voluntaria. La persona participante puede
decidir no responder determinadas preguntas o retirarse en
cualquier momento, sin consecuencias académicas ni personales.

\subsection*{A.4. Confidencialidad y tratamiento de datos}

La información obtenida se utilizará exclusivamente con fines
académicos relacionados con la actividad.

Los documentos firmados y las grabaciones se conservarán
en un espacio institucional privado, con acceso restringido
a las personas autorizadas para su revisión.

No se publicarán datos personales, grabaciones ni formularios
firmados en el repositorio público.

\subsection*{A.5. Grabación de audio}

Se solicita autorización para registrar la entrevista
únicamente mediante audio. No se realizará grabación de video.

La autorización para grabar es independiente de la decisión
de participar. Si no se autoriza, no se realizará la grabación.

\subsection*{A.6. Declaración de consentimiento}

Declaro que he recibido información sobre el propósito,
procedimiento y condiciones de la actividad, y que he podido
resolver mis dudas antes de tomar una decisión.

\vspace{0.3cm}

\noindent\textbf{Participación:}\\[0.15cm]
$\square$ Acepto participar \qquad
$\square$ No acepto participar.

\vspace{0.3cm}

\noindent\textbf{Grabación de audio:}\\[0.15cm]
$\square$ Autorizo \qquad
$\square$ No autorizo.

\vspace{0.5cm}

\noindent\textbf{Nombre del participante:}
\rule{7cm}{0.4pt}

\vspace{0.45cm}

\noindent\textbf{Firma del participante:}
\rule{7.3cm}{0.4pt}

\vspace{0.45cm}

\noindent\textbf{Fecha:}
\rule{5cm}{0.4pt}

\vspace{0.45cm}

\noindent\textbf{Nombre del entrevistador:}
\rule{7cm}{0.4pt}

\vspace{0.45cm}

\noindent\textbf{Firma del entrevistador:}
\rule{7.3cm}{0.4pt}

\vfill

\begin{center}
\footnotesize\itshape
Esta plantilla se publica sin completar.
El formulario firmado, si se utiliza, permanecerá
en el almacenamiento institucional privado.
\end{center}

\end{document}